% Bagian: sets-functions-relations
% Bab: relations
% Subbagian: equivalence-relations

\documentclass[../../../include/open-logic-section]{subfiles}

\begin{document}

\olfileid[id]{sfr}{rel}{eqv}

\olsection{Relasi Ekuivalensi}

Relasi identitas pada suatu himpunan bersifat refleksif, simetris, dan
transitif. Relasi~$R$ yang memiliki ketiga sifat tersebut sangat umum.

\begin{defn}[Relasi ekuivalensi]
Suatu relasi $R \subseteq A^2$ yang bersifat refleksif, simetris, dan transitif
disebut \emph{relasi ekuivalensi}. !!^{element}s $x$ dan $y$ dari~$A$ disebut
\emph{ekuivalen terhadap~$R$} jika~$Rxy$.
\end{defn}

Relasi ekuivalensi melahirkan gagasan \emph{kelas ekuivalensi}. Suatu relasi
ekuivalensi ``memecah'' domain menjadi blok-blok partisi yang berbeda. Di dalam
setiap blok, semua objek saling berelasi; tidak ada objek dari blok yang
berbeda yang saling berelasi. Kadang-kadang akan membantu jika kita dapat
membicarakan blok-blok ini \emph{secara langsung}. Untuk itu, kita perkenalkan
definisi berikut.

\begin{defn}\ollabel{def:equivalenceclass}
Misalkan $R \subseteq A^2$ merupakan relasi ekuivalensi. Untuk setiap $x \in
A$, \emph{kelas ekuivalensi} dari $x$ dalam~$A$ adalah himpunan
$\equivrep{x}{R} = \Setabs{y \in A}{Rxy}$. \emph{Hasil bagi} $A$ terhadap~$R$
adalah $\equivclass{A}{R} = \Setabs{\equivrep{x}{R}}{x \in A}$, yakni himpunan
semua kelas ekuivalensi tersebut.
\end{defn}

Hasil berikut membenarkan definisi kelas ekuivalensi dengan membuktikan bahwa
kelas-kelas ekuivalensi memang merupakan blok-blok partisi dari~$A$.

\begin{prop}
Jika $R \subseteq A^2$ merupakan relasi ekuivalensi, maka $Rxy$ jika dan hanya
jika $\equivrep{x}{R} = \equivrep{y}{R}$.
\end{prop}

\begin{proof}
Untuk arah kiri ke kanan, andaikan $Rxy$, dan misalkan $z \in
\equivrep{x}{R}$. Menurut definisi, $Rxz$. Karena $R$ merupakan relasi
ekuivalensi, $Ryz$. (Secara terperinci: karena $Rxy$ dan~$R$ simetris, kita
memperoleh $Ryx$; karena $Rxz$ dan~$R$ transitif, kita memperoleh~$Ryz$.) Jadi,
$z \in \equivrep{y}{R}$. Dengan menggeneralisasi,
$\equivrep{x}{R} \subseteq \equivrep{y}{R}$. Akan tetapi, dengan penalaran yang
persis sama, $\equivrep{y}{R} \subseteq \equivrep{x}{R}$. Jadi,
$\equivrep{x}{R} = \equivrep{y}{R}$ berdasarkan ekstensionalitas.

Untuk arah kanan ke kiri, andaikan $\equivrep{x}{R} = \equivrep{y}{R}$. Karena
$R$ refleksif, $Ryy$, sehingga $y \in \equivrep{y}{R}$. Jadi, berlaku pula
$y \in \equivrep{x}{R}$ berdasarkan asumsi bahwa
$\equivrep{x}{R} = \equivrep{y}{R}$.
Jadi, $Rxy$.
\end{proof}

\begin{ex}
Contoh yang baik untuk relasi ekuivalensi berasal dari aritmetika modular.
Untuk sembarang bilangan asli $a$, $b$, dan bilangan bulat positif
$n \in \PosInt$, tetapkan bahwa $a \equiv_n b$
jika dan hanya jika pembagian $a$ dengan~$n$ menghasilkan sisa yang sama
dengan pembagian $b$ dengan~$n$. (Dalam bentuk yang lebih simbolis:
$a \equiv_n b$ jika dan hanya jika, untuk suatu $k \in \Int$, $a - b = kn$.)
Sekarang, $\equiv_n$ merupakan relasi ekuivalensi untuk setiap~$n$. Tepat ada
$n$ kelas ekuivalensi berbeda yang dihasilkan oleh~$\equiv_n$; dengan kata
lain, $\equivclass{\Nat}{\equiv_n}$ mempunyai $n$ !!{element}s. Kelas-kelas
itu adalah: himpunan bilangan yang habis dibagi $n$, yakni
$\equivrep{0}{\equiv_n}$; himpunan bilangan yang jika dibagi $n$ bersisa~$1$,
yakni $\equivrep{1}{\equiv_n}$; \ldots; dan himpunan bilangan yang jika dibagi
$n$ bersisa~$n-1$, yakni~$\equivrep{n-1}{\equiv_n}$.
\end{ex}

\begin{prob}
Tunjukkan bahwa $\equiv_n$ merupakan relasi ekuivalensi untuk setiap $n \in
\PosInt$, dan bahwa $\equivclass{\Nat}{\equiv_n}$ mempunyai tepat $n$ anggota.
\end{prob}

\end{document}
